\noindent\textbf{Resultado.} $y=2(1+Ce^{4x})/(1-Ce^{4x})$ donde $1-Ce^{4x}\ne0$, junto con $y=-2$. El caso $C=0$ ya da $y=2$.

\noindent\textit{Desarrollo.} Los equilibrios satisfacen $y^2-4=0$: son $y=\pm2$. Fuera de ellos, separar e integrar da $\tfrac14\ln|(y-2)/(y+2)|=x+K$. Al absorber el signo en una constante real se obtiene $(y-2)/(y+2)=Ce^{4x}$ y la familia indicada. Si $C<0$, el denominador es positivo y la solución está definida en toda $\mathbb R$, entre $-2$ y $2$. Si $C>0$, hay un polo en $x_*=-\ln C/4$ y los intervalos máximos son $(-\infty,x_*)$ y $(x_*,\infty)$. Si $C=0$, la solución es el equilibrio $2$, y el equilibrio $-2$ se añade por separado. Con $z=Ce^{4x}$, la derivada de la familia es $16z/(1-z)^2$, y también $y^2-4=16z/(1-z)^2$. El campo es suave; no existen empalmes de soluciones distintas en puntos finitos de los equilibrios.
